\documentstyle[12pt]{article}
\pagestyle{empty}

\begin{document}

{
\baselineskip10mm
\centerline{\LARGE\bf Applications of Computer Algebra}
\centerline{\LARGE\bf in}
\centerline{\LARGE\bf Scientific Computing}
}

\vspace{1.0cm}

\begin{center}
{\small
\begin{tabular}{c}
	{\normalsize Hans-Joachim Bungartz} \\
	{FORTWIHR --- The Bavarian Consortium} \\
        {for High Performance Scientific Computing} \\
	{c/o Institut f\"ur Informatik der}\\
	{Technischen Universit\"at M\"unchen }\\
	{D-80290 M\"unchen, Germany}\\
        {\tt\footnotesize bungartz@informatik.tu-muenchen.de}
\end{tabular}
}
\end{center}

\vspace{1cm}
\centerline{{\large\bf Abstract}}
\vspace{1cm}

In the field of scientific computing, the methods and tools of classical
numerical analysis are traditionally much more common than the use of
computer algebra systems. Nevertheless, modern computer algebra systems
can support both the development, the theoretical analysis, and the 
efficient use of algorithms and codes in the different areas of application
of scientific computing. Therefore, it turns out to be more and more
important to bring together the two worlds of numerical and
symbolic computation.

We report a couple of situations where computer algebra tools, recently,
have been successfully integrated in scientific computing projects. Here,
the areas of application range from support in proving approximation 
properties of discretization schemes to data compression for purposes
of visualization and multi media applications. Finally, results from
the implementation of computer algebra modules in programs for 
computational fluid dynamics are presented.

\end{document}
